% !TEX program = pdflatex

\documentclass[a4paper,oneside,14pt]{extreport}

% --- Настройка кодировки и языков ---
\usepackage[T1,T2A]{fontenc}
\usepackage[utf8]{inputenc}
\usepackage[english,main=russian]{babel}
\usepackage{fix-cm}

% --- Настройка полей ---
\usepackage[
	left=30mm,
	right=10mm,
	top=20mm,
	bottom=20mm,
]{geometry}

% --- Настройка переносов ---
\usepackage{microtype}
\sloppy

% --- Настройка межстрочного интервала ---
\usepackage{setspace}
\onehalfspacing

% --- Настройка абзацного отступа ---
\usepackage{indentfirst}
\setlength{\parindent}{12.5mm}

% --- Настройка заголовков ---
\makeatletter
\renewcommand\LARGE{\@setfontsize\LARGE{22pt}{20}}
\renewcommand\Large{\@setfontsize\Large{20pt}{20}}
\renewcommand\large{\@setfontsize\large{16pt}{20}}
\makeatother
\usepackage{titlesec}
\titleformat{\chapter}[block]{\hspace{\parindent}\large\bfseries}{\thechapter}{0.5em}{\large\bfseries\raggedright}
\titleformat{name=\chapter,numberless}[block]{\hspace{\parindent}}{}{0pt}{\large\bfseries\centering}
\titleformat{\section}[block]{\hspace{\parindent}\large\bfseries}{\thesection}{0.5em}{\large\bfseries\raggedright}
\titleformat{\subsection}[block]{\hspace{\parindent}\large\bfseries}{\thesubsection}{0.5em}{\large\bfseries\raggedright}
\titleformat{\subsubsection}[block]{\hspace{\parindent}\large\bfseries}{\thesubsection}{0.5em}{\large\bfseries\raggedright}
\titlespacing{\chapter}{12.5mm}{-22pt}{10pt}
\titlespacing{\section}{12.5mm}{10pt}{10pt}
\titlespacing{\subsection}{12.5mm}{10pt}{10pt}
\titlespacing{\subsubsection}{12.5mm}{10pt}{10pt}

% --- Пакеты ---
\usepackage{xcolor}
\usepackage{tabularx}
\usepackage{longtable}
\usepackage{booktabs}
\usepackage[labelsep=endash,figurename=Рисунок,singlelinecheck=false]{caption}
\captionsetup[figure]{justification=centering}
\usepackage{lscape}
\usepackage{afterpage}
\usepackage{amsmath}
\usepackage{amssymb}
\usepackage{csquotes}
\usepackage[unicode,hidelinks]{hyperref}
\usepackage{xifthen}
\usepackage{graphicx}
\usepackage{float}
\usepackage{wrapfig}
\usepackage{tikzscale}
\usepackage{pgfplots}
\pgfplotsset{compat=newest}

% --- Титульный лист (bmstu-title.sty) ---
\newcommand{\documentmeta}[4]{
	\hypersetup{
		pdftitle={#1 #2 #3},
		pdfsubject={#4},
		pdfauthor={#2}
	}
}
\usepackage[normalem]{ulem}
\usepackage{stackengine}
\newcommand{\longunderline}[1]{\uline{#1\hfill\mbox{}}}
\newcommand{\fixunderline}[3]{$\underset{\text{#3}}{\text{\uline{\stackengine{0pt}{\hspace{#2}}{\text{#1}}{O}{c}{F}{F}{L}}}}$}
\makeatletter
\newcommand{\vhrulefill}[1]{\leavevmode\leaders\hrule\@height#1\hfill \kern\z@}
\makeatother

\newcommand{\titlepageheader}[2]{
	\begin{wrapfigure}[7]{l}{0.14\linewidth}
		\vspace{3.4mm}
		\hspace{-8mm}
		\includegraphics[width=0.89\linewidth]{bmstu-logo}
	\end{wrapfigure}
	{
		\singlespacing \small
		Министерство науки и высшего образования Российской Федерации \\
		Федеральное государственное бюджетное образовательное учреждение \\
		высшего образования \\
		<<Московский государственный технический университет \\
		имени Н.~Э.~Баумана \\
		(национальный исследовательский университет)>> \\
		(МГТУ им. Н.~Э.~Баумана) \\
	}
	\vspace{-4.2mm}
	\vhrulefill{0.9mm} \\
	\vspace{-7mm}
	\vhrulefill{0.2mm} \\
	\vspace{2.8mm}
	{
		\small
		ФАКУЛЬТЕТ \longunderline{<<#1>>} \\
		\vspace{3.3mm}
		КАФЕДРА \longunderline{<<#2>>} \\
	}
}

\newcommand{\titlepageresearchtitle}[1]{
	{
		\LARGE \bfseries
		РАСЧЕТНО-ПОЯСНИТЕЛЬНАЯ ЗАПИСКА \\
	}
	\vspace{5mm}
	{
		\Large \itshape
		К НАУЧНО-ИССЛЕДОВАТЕЛЬСКОЙ РАБОТЕ \\
		\vspace{5mm}
		НА ТЕМУ: \\
		<<#1>> \\
	}
}

\usepackage{pgffor}
\newcommand*\titlepagestudentscontent{}
\newcommand{\maketitlepagestudent}[1]{
	\foreach \s/\g in {#1} {
		\gappto\titlepagestudentscontent{Студент \fixunderline}
		\xappto\titlepagestudentscontent{{\g}}
		\gappto\titlepagestudentscontent{{25mm}{(Группа)} &}
		\gappto\titlepagestudentscontent{\fixunderline{}{40mm}{(Подпись, дата)} \vspace{1.3mm} &}
		\gappto\titlepagestudentscontent{\fixunderline}
		\xappto\titlepagestudentscontent{{\s}}
		\gappto\titlepagestudentscontent{{40mm}{(И.~О.~Фамилия)} \\}
	}
}

\newcommand*\titlepageotherscontent{}
\newcommand{\maketitlepageothers}[2]{
	\foreach \c in {#2} {
		\gappto\titlepagestudentscontent{#1 &}
		\gappto\titlepagestudentscontent{\fixunderline{}{40mm}{(Подпись, дата)} \vspace{1.3mm} &}
		\gappto\titlepagestudentscontent{\fixunderline}
		\xappto\titlepagestudentscontent{{\c}}
		\gappto\titlepagestudentscontent{{40mm}{(И.~О.~Фамилия)} \\}
	}
}

\newcommand{\titlepageauthors}[5]{
	{
		\renewcommand{\titlepagestudentscontent}{}
		\maketitlepagestudent{#1}
		\renewcommand{\titlepageotherscontent}{}
		\maketitlepageothers{#2}{#3}
		\maketitlepageothers{Консультант}{#4}
		\maketitlepageothers{Нормоконтролер}{#5}
		\small
		\begin{tabularx}{\textwidth}{@{}>{\hsize=.5\hsize}X>{\hsize=.25\hsize}X>{\hsize=.25\hsize}X@{}}
			\titlepagestudentscontent
			\titlepageotherscontent
		\end{tabularx}
	}
}

\newcommand{\makeresearchtitle}[6]{
	\documentmeta{Отчет по НИР}{#4}{}{#3}
	\begin{titlepage}
		\centering
		\titlepageheader{#1}{#2}
		\vspace{15.8mm}
		\titlepageresearchtitle{#3}
		\vfill
		\titlepageauthors{#4}{Руководитель НИР}{#5}{#6}{}
		\vspace{14mm}
		\textit{{\the\year} г.}
	\end{titlepage}
	\setcounter{page}{2}
}

% --- Реферат (bmstu-essay.sty) ---
\usepackage{lastpage}
\usepackage{etoolbox}
\newcounter{totfigures}
\newcounter{tottables}
\providecommand\totfig{}
\providecommand\tottab{}
\makeatletter
\AtEndDocument{%
	\addtocounter{totfigures}{\value{figure}}%
	\addtocounter{tottables}{\value{table}}%
	\immediate\write\@mainaux{%
		\string\gdef\string\totfig{\number\value{totfigures}}%
		\string\gdef\string\tottab{\number\value{tottables}}%
	}%
}
\makeatother
\pretocmd{\chapter}{\addtocounter{totfigures}{\value{figure}}\setcounter{figure}{0}}{}{}
\pretocmd{\chapter}{\addtocounter{tottables}{\value{table}}\setcounter{table}{0}}{}{}



\newenvironment{essay}[1]{
	\chapter*{РЕФЕРАТ}
	\addcontentsline{toc}{chapter}{РЕФЕРАТ}
	Расчетно-пояснительная записка \begin{NoHyper}\pageref{LastPage}\end{NoHyper}~с., \totfig~рис., \tottab~табл., 10~источн., 1~прил.
	
	\vspace{1em}
	\noindent \MakeUppercase{#1} \par
}{}

% --- Обозначения и сокращения (bmstu-defabbr.sty) ---
\usepackage{enumitem}
\newcommand{\definition}[2]{\item \noindent #1 --- #2}
\newenvironment{abbreviations}{
	\chapter*{ОБОЗНАЧЕНИЯ И СОКРАЩЕНИЯ}
	\addcontentsline{toc}{chapter}{ОБОЗНАЧЕНИЯ И СОКРАЩЕНИЯ}
	В настоящей расчетно-пояснительной записке применяют следующие сокращения и обозначения.
	\begin{description}[leftmargin=0pt]
}{
	\end{description}
}

% --- Оглавление (bmstu-toc.sty) ---
\newcommand{\maketableofcontents}{
	\renewcommand\contentsname{СОДЕРЖАНИЕ}
	\tableofcontents
}

% --- Приложения ---
\usepackage[titletoc]{appendix}

% --- Настройка списков ГОСТ ---
\setlist[itemize]{label=---, wide, labelindent=\parindent, labelsep=0.5em, itemsep=2pt, parsep=0pt, topsep=4pt}

\makeatletter
\AddEnumerateCounter{\asbuk}{\russian@alph}{щ}
\makeatother
\setlist[enumerate,1]{label=\arabic*), wide, labelindent=\parindent, labelsep=0.5em, itemsep=2pt, parsep=0pt, topsep=4pt}
\setlist[enumerate,2]{label=\asbuk*), wide, labelindent=\parindent, labelsep=0.5em, itemsep=2pt, parsep=0pt, topsep=4pt}

\newcommand{\blank}[1][1.5cm]{\underline{\hspace{#1}}}

\begin{document}

% ======================================================================
%  ТИТУЛЬНЫЙ ЛИСТ
% ======================================================================
\makeresearchtitle
{Информатика и системы управления (ИУ)} % факультет
{Информационные системы и телекоммуникации (ИУ3)} % кафедра
{Проектирование структурированной кабельной системы офиса лаборатории тестирования устройств IoT} % тема
{А.\,А.\,Захватов/ИУ3-33Б} % студент / группа
{И.\,М.\,Сидякин} % руководитель
{} % консультант

% ======================================================================
%  РЕФЕРАТ
% ======================================================================
\begin{essay}{структурированная кабельная система, иерархическая звезда, проектирование СКС, телекоммуникационный шкаф, видеонаблюдение, активное сетевое оборудование, маршрутизация, импортозамещение}

Объектом разработки является структурированная кабельная система (СКС) этажа лаборатории тестирования устройств IoT.

Цель работы~--- разработка СКС этажа по архитектуре иерархической звезды с одной серверной и одной кроссовой комнатой.

В~процессе работы были выполнены расчеты количества информационных розеток, емкости коммутационных панелей и плотности заполнения телекоммуникационных шкафов. Проведен расчет трафика камер видеонаблюдения и нагрузки на медиасервер и операторскую в~штатном (MCU) и аварийном (SFU) режимах. Обоснована пропускная способность каналов связи и осуществлен выбор камер с~поддержкой кодеков H.265/H.265+.

Спроектирована физическая инфраструктура сети: рассчитаны контрольные длины кабельных линий, сечения и уровни заполнения кабельных каналов и лотков, а~также механические нагрузки на элементы крепления. Выполнен тепловой расчет, включающий оценку тепловыделения, параметров охлаждения и требуемой емкости ИБП для шкафов серверной и кроссовой. Разработана и обоснована компоновка телекоммуникационного оборудования с~учетом требований к~заземлению.

В~результате выполнения работы разработано техническое задание на СКС в~соответствии с~требованиями ГОСТ~34.602-2020 и ГОСТ~Р~15.301-2016. Осуществлен подбор активного сетевого оборудования из Единого реестра российской радиоэлектронной продукции Минпромторга~РФ и~произведен проверочный расчет спроектированной вычислительной сети. Подготовлен полный комплект отчетной документации, включающий расчетно-пояснительную записку, чертеж схемы сети объекта на~листе формата~А2 по~стандартам ЕСКД, схему логической топологии и ведомость покупных изделий.
\end{essay}

% ======================================================================
%  ОПРЕДЕЛЕНИЯ, ОБОЗНАЧЕНИЯ И СОКРАЩЕНИЯ
% ======================================================================
\begin{abbreviations}
\definition{АКБ}{аккумуляторная батарея}
\definition{БП}{блок питания}
\definition{БРП}{блок распределения питания}
\definition{ГОСТ~Р}{государственный стандарт России}
\definition{ЕСКД}{единая система конструкторской документации}
\definition{ИБП}{источник бесперебойного питания}
\definition{КНС}{кабеленесущие системы}
\definition{КС}{кабельная система}
\definition{КПД}{коэффициент полезного действия}
\definition{МГТУ~им.~Н.\,Э.~Баумана}{Федеральное государственное автономное образовательное учреждение высшего образования <<Московский государственный технический университет имени Н.\,Э.~Баумана (национальный исследовательский университет)>>}
\definition{НИР}{научно-исследовательская работа}
\definition{РПЗ}{расчетно-пояснительная записка}
\definition{РФ}{Российская Федерация}
\definition{СКС}{структурированная кабельная система}
\definition{СКУД}{система контроля и управления доступом}
\definition{ТЗ}{техническое задание}
\definition{ТО}{техническое обслуживание}
\definition{ТС}{технические средства кабельной системы}
\definition{ЭРП}{этажный распределительный пункт}
\definition{IoT}{Internet of Things (интернет вещей)}
\definition{MCU}{Multipoint Control Unit (центральный узел микширования потоков)}
\definition{MUTOA}{Multi-User Telecommunications Outlet Assembly (многопользовательская телекоммуникационная розетка)}
\definition{PoE}{Power over Ethernet (питание по кабелю Ethernet)}
\definition{RJ-45}{Registered Jack~45}
\definition{SFU}{Selective Forwarding Unit (узел селективной ретрансляции потоков)}
\definition{U}{Unit}
\end{abbreviations}

% ======================================================================
%  СОДЕРЖАНИЕ
% ======================================================================
\maketableofcontents

% ======================================================================
%  ВВЕДЕНИЕ
% ======================================================================
\chapter*{ВВЕДЕНИЕ}
\addcontentsline{toc}{chapter}{ВВЕДЕНИЕ}

В~условиях стремительного развития цифровых технологий и роста объемов передаваемых данных надежная телекоммуникационная инфраструктура является критически важным элементом любого современного предприятия. Особые требования к~производительности, отказоустойчивости и безопасности сетей предъявляются в~специализированных испытательных центрах и лабораториях, где осуществляется работа с~большими массивами тестовых данных и непрерывный мониторинг процессов.

Актуальность разработки структурированной кабельной системы (СКС) для лаборатории тестирования устройств IoT обусловлена необходимостью создания единой, масштабируемой и высокоскоростной физической среды передачи сигналов. Современная СКС должна не~только обеспечивать информационный обмен между лабораторными стендами, операторскими рабочими местами и серверами, но~и~поддерживать интеграцию систем физической безопасности: видеонаблюдения высокого разрешения и~систем контроля и~управления доступом (СКУД) с~подачей электропитания по~технологии PoE+/PoE++. Дополнительную актуальность проекту придает необходимость реализации принципов импортозамещения при~выборе активного сетевого оборудования.

Объектом исследования и разработки является телекоммуникационная инфраструктура этажа лаборатории тестирования устройств IoT.

Предметом разработки выступают методы проектирования физического уровня локальных вычислительных сетей, принципы расчета кабельных трасс, энергетического баланса и параметров микроклимата коммутационных узлов.

Целью данной научно-исследовательской работы является проектирование масштабируемой и надежной структурированной кабельной системы этажа лаборатории на~базе архитектуры иерархической звезды, отвечающей современным стандартам пропускной способности и отказоустойчивости.

Для достижения поставленной цели в~рамках работы сформулированы и решены следующие задачи:
\begin{itemize}
    \item проведен анализ архитектурного плана и исходных данных объекта, осуществлено оптимальное размещение лабораторного, операторского и телекоммуникационного оборудования;
    \item выполнен расчет необходимого количества информационных портов, емкости коммутационных панелей и уровня заполнения телекоммуникационных шкафов в~серверной и кроссовой;
    \item произведен расчет пропускной способности каналов связи с~учетом трафика камер видеонаблюдения в~штатном (MCU) и аварийном (SFU) режимах;
    \item спроектированы кабельные трассы: рассчитаны контрольные длины линий, заполняемость кабельных каналов и лотков, а~также механические нагрузки на~элементы крепления;
    \item выполнен комплексный расчет параметров микроклимата и энергоснабжения, включающий оценку тепловыделения, выбор систем охлаждения и расчет требуемой емкости источников бесперебойного питания (ИБП);
    \item осуществлен подбор активного сетевого оборудования из~Единого реестра российской радиоэлектронной продукции Минпромторга~РФ с~последующим проверочным расчетом сети;
    \item разработано техническое задание на~СКС и подготовлен полный комплект проектной и конструкторской документации.
\end{itemize}


% ======================================================================
%  РАСЧЁТ ПРОПУСКНОЙ СПОСОБНОСТИ КАНАЛОВ СВЯЗИ
% ======================================================================

\chapter{Расчёт пропускной способности каналов связи}

\section{Общее число портов СКС}

Общее число портов СКС определяется как сумма портов по всем категориям помещений:
\begin{equation}
    N_{\text{общ}} = N_{\text{лаб}} + N_{\text{опер}} + N_{\text{проч}},
    \label{eq:n_total}
\end{equation}
где:
\begin{itemize}
    \item $N_{\text{лаб}}$~--- порты в~лабораториях;
    \item $N_{\text{опер}}$~--- порты в~операторской;
    \item $N_{\text{проч}}$~--- порты в~прочих помещениях.
\end{itemize}

\textbf{Порты в~лабораториях.} Согласно варианту~V10, под лаборатории отведено 2~помещения (№\,4 и~№\,11), в~которых размещено 6~лабораторных стендов (по 3~стенда на~лабораторию). В~каждом стенде устанавливается многопользовательская розетка (MUTOA), ёмкость которой определяется числом тестируемых IoT-устройств, портом узконаправленной камеры контроля и резервными портами для контрольно-измерительного оборудования (КИП).

Для помещения №\,4 (Лаб.\,1) ёмкость MUTOA одного стенда составляет:
\begin{equation}
    N_{\text{MUTOA\_лаб1}} = N_{\text{IoT}} + N_{\text{КИП}} + N_{\text{кам}} = 18 + 6 + 1 = 25 \text{ портов}.
    \label{eq:mutoa_lab1}
\end{equation}

Для помещения №\,11 (Лаб.\,2) ёмкость MUTOA одного стенда составляет:
\begin{equation}
    N_{\text{MUTOA\_лаб2}} = N_{\text{IoT}} + N_{\text{КИП}} + N_{\text{кам}} = 8 + 6 + 1 = 15 \text{ портов}.
    \label{eq:mutoa_lab2}
\end{equation}

Суммарная ёмкость MUTOA лабораторных стендов для помещения №\,4:
\begin{equation}
    N_{\text{лаб1\_MUTOA}} = 3 \times 25 = 75 \text{ портов}.
\end{equation}

Суммарная ёмкость MUTOA лабораторных стендов для помещения №\,11:
\begin{equation}
    N_{\text{лаб2\_MUTOA}} = 3 \times 15 = 45 \text{ портов}.
\end{equation}

Дополнительно в~каждой из 2~лабораторий размещаются IP-камера общего вида и контроллер СКУД, подключаемые к~портам PoE+. Это добавляет:
\begin{equation}
    N_{\text{лаб\_доп}} = 2 \times (1_{\text{камера}} + 1_{\text{СКУД}}) = 4 \text{ порта}.
\end{equation}

Итого портов в~лабораториях:
\begin{equation}
    N_{\text{лаб}} = 75 + 2 + 45 + 2 = 124 \text{ порта}.
\end{equation}

(Проверка: помещение №\,4: $3 \times 25 + 1 + 1 = 77$ портов; помещение №\,11: $3 \times 15 + 1 + 1 = 47$ портов; $77 + 47 = 124$ порта.)

\textbf{Порты в~операторской.} Операторская размещается в~помещении №\,20 (<<Опенспейс>>) и включает 6~рабочих мест (постов). На~каждом посту размещается 1~моноблок (1~порт PoE++), 1~рабочая станция (1~порт) и 1~резервная розетка. Также в~общей зоне операторской устанавливаются 2~дополнительные розетки:
\begin{equation}
    N_{\text{опер}} = 6 \times (1_{\text{мон}} + 1_{\text{РС}} + 1_{\text{рез}}) + 2_{\text{общ}} = 18 + 2 = 20 \text{ портов}.
    \label{eq:n_oper}
\end{equation}

\textbf{Порты в~прочих помещениях.} Двумя розетками СКС оборудуются 3~помещения (№\,7, №\,18, №\,24):
\begin{equation}
    N_{\text{проч}} = 3 \times 2 = 6 \text{ портов}.
    \label{eq:n_other}
\end{equation}

Итого общее число портов СКС:
\begin{equation}
    N_{\text{общ}} = 124 + 20 + 6 = 150 \text{ портов}.
    \label{eq:n_total_calc}
\end{equation}

\textbf{Разграничение по типам PoE:}
\begin{itemize}
    \item PoE++ (IEEE~802.3bt, до~90~Вт): моноблоки в~операторской~--- 6~портов;
    \item PoE+ (IEEE~802.3at, до~30~Вт): камеры контроля стендов (6~портов) + камеры общего вида (2~порта) + контроллеры СКУД (2~порта)~--- 10~портов;
    \item Без PoE: IoT-устройства, резервные порты MUTOA, рабочие станции и резервные розетки операторской, розетки прочих помещений~--- 134~порта.
\end{itemize}

Итого: PoE++~--- 6~портов, PoE+~--- 10~портов, без~PoE~--- 134~порта.

% ---- Расчёт нагрузки от видео трафика ----
\section{Расчёт нагрузки от видео трафика}

Согласно техническому заданию, во время тестов в~лабораториях нагрузка формируется в~основном за~счёт видео трафика, к~которому добавляется незначительный по~объёму поток команд и телеметрических данных.

Формула для расчёта видеотрафика без применения кодека (несжатый поток):
\begin{equation}
    P' = \frac{W \cdot H \cdot FPS \cdot Bit}{10^{6}} \quad [\text{Мбит/с}],
    \label{eq:bitrate_raw}
\end{equation}
где:
\begin{itemize}
    \item $W$~--- ширина изображения в~пикселях;
    \item $H$~--- высота изображения в~пикселях;
    \item $FPS$~--- количество кадров в~секунду;
    \item $Bit$~--- количество бит для кодирования одного пикселя (при цветовой модели YCbCr~4:2:2 принимается $Bit = 16$).
\end{itemize}

Расчёт несжатых потоков для применяемых камер:
\begin{equation}
    P'_{\text{узк\_MCU}} = \frac{3840 \cdot 2160 \cdot 60 \cdot 16}{10^{6}} = 7962{,}6 \text{ Мбит/с},
    \label{eq:raw_narrow}
\end{equation}
\begin{equation}
    P'_{\text{общ\_MCU}} = \frac{3840 \cdot 2160 \cdot 30 \cdot 16}{10^{6}} = 3981{,}3 \text{ Мбит/с},
    \label{eq:raw_overview_mcu}
\end{equation}
\begin{equation}
    P'_{\text{общ\_SFU}} = \frac{640 \cdot 480 \cdot 25 \cdot 16}{10^{6}} = 122{,}88 \text{ Мбит/с}.
    \label{eq:raw_overview_sfu}
\end{equation}

При использовании сжатия кодеком H.265~/~H.265+ битрейт видеопотоков составляет:
\begin{itemize}
    \item $P_{\text{узк\_MCU}} = 18$~Мбит/с~--- нагрузка на узконаправленную камеру (4K 60~Гц) в~режиме MCU;
    \item $P_{\text{общ\_MCU}} = 9$~Мбит/с~--- нагрузка на камеру общего вида (4K 25--30~Гц) в~режиме MCU;
    \item $P_{\text{общ\_SFU}} = 1$~Мбит/с~--- нагрузка на камеру общего вида (640$\times$480 25~Гц) в~режиме SFU.
\end{itemize}

% ---- Расчёт нагрузки на коммутаторы доступа ----
\section{Расчёт нагрузки на коммутаторы доступа}

В~штатном режиме MCU на каждую лабораторию приходится нагрузка от 3~узконаправленных камер и 1~обзорной камеры:
\begin{equation}
    P_{\text{лаб1\_MCU}} = 3 \times P_{\text{узк\_MCU}} + 1 \times P_{\text{общ\_MCU}} = 3 \times 18 + 9 = 63 \text{ Мбит/с}.
    \label{eq:p_lab1}
\end{equation}
\begin{equation}
    P_{\text{лаб2\_MCU}} = 3 \times P_{\text{узк\_MCU}} + 1 \times P_{\text{общ\_MCU}} = 3 \times 18 + 9 = 63 \text{ Мбит/с}.
    \label{eq:p_lab2}
\end{equation}

Суммарный видеотрафик от 2~лабораторий, поступающий на коммутаторы доступа:
\begin{equation}
    P_{\text{доступа\_MCU}} = 63 + 63 = 126 \text{ Мбит/с}.
    \label{eq:p_access_mcu}
\end{equation}

В~аварийном режиме SFU направленные камеры отключаются, а~обзорные камеры переключаются в~режим 640$\times$480:
\begin{equation}
    P_{\text{доступа\_SFU}} = 2 \times P_{\text{общ\_SFU}} = 2 \times 1 = 2 \text{ Мбит/с}.
    \label{eq:p_access_sfu}
\end{equation}

Проверка: 126~Мбит/с $<$ 1~Гбит/с (производительность портов 1G обеспечена с~большим запасом).

% ---- Расчёт нагрузки на медиасервер ----
\section{Расчёт нагрузки на медиасервер}

\textbf{Штатный режим MCU:}
\begin{itemize}
    \item Входящий поток от лабораторий: $P_{\text{медиа\_вх}} = 126$~Мбит/с;
    \item Исходящий поток на 6~моноблоков (4K 60~Гц): $P_{\text{медиа\_исх}} = 6 \times 18 = 108$~Мбит/с;
    \item Суммарная нагрузка на интерфейс медиасервера: $P_{\text{медиа\_MCU}} = 126 + 108 = 234$~Мбит/с.
\end{itemize}

\textbf{Аварийный режим SFU:}
\begin{itemize}
    \item Входящий поток от обзорных камер: $P_{\text{медиа\_вх}} = 2$~Мбит/с;
    \item Исходящий поток (трансляция 2~обзорных камер на 6~моноблоков): $P_{\text{медиа\_исх}} = 6 \times 2 \times 1 = 12$~Мбит/с;
    \item Суммарная нагрузка: $P_{\text{медиа\_SFU}} = 2 + 12 = 14$~Мбит/с.
\end{itemize}

% ---- Расчёт нагрузки на коммутатор ядра ----
\section{Расчёт нагрузки на коммутатор ядра}

Всего тестируемых IoT-устройств: $N_{\text{IoT}} = 78$~шт. ($3 \times 18 + 3 \times 8$). Скорость потока при прошивке 1~устройства~--- 50~Мбит/с. Пиковый трафик прошивок:
\begin{equation}
    P_{\text{IoT\_тест}} = 78 \times 50 = 3900 \text{ Мбит/с}.
    \label{eq:p_iot}
\end{equation}

Нагрузка при восходящем соединении (от доступа к~серверному сегменту):
\begin{equation}
    P_{\text{uplink}} = P_{\text{доступа\_MCU}} + P_{\text{IoT\_тест}} + P_{\text{CI/CD/NAS}} = 126 + 3900 + 1000 = 5026 \text{ Мбит/с}.
    \label{eq:p_uplink}
\end{equation}

Нагрузка при нисходящем соединении (от медиасервера и NAS к~операторской и рабочим станциям):
\begin{equation}
    P_{\text{downlink}} = P_{\text{медиа\_исх}} + P_{\text{доступа\_MCU}} + P_{\text{данные}} = 108 + 126 + 1200 = 1434 \text{ Мбит/с}.
    \label{eq:p_downlink}
\end{equation}

Итоговая нагрузка на коммутатор ядра:
\begin{equation}
    P_{\text{ядра}} = \max(5026,\; 1434) = 5026 \text{ Мбит/с} \approx 5{,}03 \text{ Гбит/с}.
    \label{eq:p_core}
\end{equation}

Проверка: 5,03~Гбит/с $>$ 1~Гбит/с, следовательно, для магистральных каналов сервера и коммутаторов требуется использование портов 10G~SFP+.

% ---- Расчёт PoE-бюджета ----
\section{Расчёт PoE-бюджета}

\begin{equation}
    P_{\text{PoE+}} = 10 \text{ портов} \times 30 \text{ Вт} = 300 \text{ Вт},
    \label{eq:poe_plus}
\end{equation}
\begin{equation}
    P_{\text{PoE++}} = 6 \text{ портов} \times 90 \text{ Вт} = 540 \text{ Вт},
    \label{eq:poe_plusplus}
\end{equation}
\begin{equation}
    P_{\text{PoE\_общ}} = 300 + 540 = 840 \text{ Вт}.
    \label{eq:poe_total}
\end{equation}

С~учётом коэффициента запаса 30\,\%:
\begin{equation}
    P_{\text{PoE\_треб}} = 840 \times 1{,}3 = 1092 \text{ Вт}.
    \label{eq:poe_required}
\end{equation}

% ---- Выбор коммутаторов доступа ----
\section{Выбор коммутаторов доступа}

Для подключения 150 абонентских портов СКС выбраны коммутаторы уровня доступа из Единого реестра российской радиоэлектронной продукции Минпромторга РФ. С целью оптимизации стоимости проекта и рационального использования монтажного пространства телекоммуникационного шкафа применена комбинированная схема распределения абонентов: устройства, не требующие дистанционного питания (всего 134 порта), частично подключаются к свободным интерфейсам коммутаторов с поддержкой PoE. Выбор оборудования, согласно данным файла image\_3cfa7c.png, разделен по требуемым стандартам питания:

\begin{enumerate}
    \item \textbf{Для устройств с питанием PoE++ (6 портов)}: выбран коммутатор Aquarius AQ-N3000-48P4Y2Q (48 портов 1G PoE++, 4 порта 25G SFP28, 2 порта 40G QSFP+, бюджет PoE 1080/2160 Вт), как указано в image\_3cfa7c.png. Достаточно 1 коммутатора. Под моноблоки операторов задействуется 6 интерфейсов, а оставшиеся 42 порта выделяются для подключения оборудования без поддержки PoE и формирования технологического резерва.
    \item \textbf{Для оборудования с питанием PoE+ (10 портов)}: выбран коммутатор Eltex MES2448P (48 портов 1G PoE+, 4 порта 10G SFP+, бюджет PoE 720 Вт). Достаточно 1 коммутатора. Под видеокамеры и сетевые контроллеры СКУД задействуется 10 интерфейсов, свободные 38 портов маршрутизируются под стандартные абонентские розетки без PoE.
    \item \textbf{Для портов без PoE (основной массив)}: выбран коммутатор Eltex MES2448B (48 портов 1G RJ-45, 4 порта 10G SFP+), на основании image\_3cfa7c.png. Из суммарно требуемых 134 портов без поддержки PoE, 80 портов уже обеспечены свободной емкостью вышеупомянутых коммутаторов Aquarius и Eltex MES2420-48P. Для покрытия оставшейся потребности (54 порта) требуется 2 коммутатора Eltex MES2448B, предоставляющих суммарно 96 портов.
\end{enumerate}

% ---- Выбор коммутатора ядра ----
\section{Выбор коммутатора ядра}

Коммутатор ядра связывает коммутаторы доступа, серверное оборудование (медиасервер, NAS-хранилище, сервер сборки CI/CD) и~главный маршрутизатор.

Суммарная нагрузка на ядро составляет 5,03~Гбит/с, что~превышает лимит одного 1G~интерфейса. В~качестве коммутатора ядра выбран Eltex~MES5324 (24~порта 10G SFP+, 2~порта 100G QSFP28), полностью удовлетворяющий требованиям по~производительности, числу портов 10G и~импортозамещению.





\subsection{Выбор главного маршрутизатора}

Для обеспечения маршрутизации трафика локальной вычислительной сети и высокоскоростного подключения к сети связи общего пользования выбран сервисный маршрутизатор Eltex ESR-1200. В соответствии с техническим заданием устройство терминирует волоконно-оптические каналы связи со скоростью не менее 1 Гбит/с. Оборудование организует основной (к главной серверной здания) и резервный (к резервной серверной) вводы от слаботочного стояка СКС, что обеспечивает отказоустойчивость внешнего канала. Наличие интерфейсов 10G SFP+ позволяет создать агрегированный канал связи с коммутатором ядра. Маршрутизатор поддерживает протоколы динамической маршрутизации, межсетевое экранирование и трансляцию сетевых адресов (NAT), что полностью удовлетворяет проектным требованиям к изоляции и маршрутизации виртуальных подсетей (VLAN) лаборатории. 

\subsection{Выбор моноблоков}

Для оснащения рабочих мест операторов в помещении опенспейса (6 постов) выбраны моноблоки «Гравитон» М75И. Максимальная заявленная потребляемая мощность данного устройства составляет 120 Вт. Поскольку стандарт IEEE 802.3bt Type 4 (PoE++) ограничивает передаваемую мощность на порту значением до 90 Вт, для обеспечения моноблоков дистанционным питанием по витой паре требуется применение специализированных Ultra PoE инжекторов (Midspan), способных выдавать необходимую мощность в 120 Вт на порт.


\subsection{Выбора видеокамер}

Система видеонаблюдения лабораторий строится на базе IP-камер с поддержкой дистанционного питания. Для детализированного контроля лабораторных стендов (6 шт.) выбраны узконаправленные IP-камеры QVC-MiR501Z1 с потребляемой мощностью 8 Вт. Для обеспечения общего контроля помещений лабораторий (2 шт.) применяются обзорные IP-камеры QTECH QVC-MiR803M с энергопотреблением 7 Вт. Обе модели камер поддерживают стандартизированное питание по технологии PoE (IEEE 802.3af/at). 



\subsection{Выбор серверного оборудования}

Вычислительная инфраструктура серверной комнаты базируется на высокопроизводительных стоечных решениях, монтируемых в 19-дюймовый телекоммуникационный шкаф. Оборудование подобрано с учетом обеспечения расчетной пропускной способности интерфейсов и резервирования аппаратных компонентов:

\begin{itemize}
    \item \textbf{Медиасервер:} Для централизованной обработки видеотрафика выбран сервер Гравитон С2242И. Устройство обладает достаточными вычислительными ресурсами для приема, перекодирования и трансляции индивидуальных 4K-видеопотоков в штатном режиме (MCU), а также для работы в аварийном режиме (SFU). Сетевая подсистема сервера включает оптические интерфейсы 10G SFP+, исключающие образование узких мест при приеме 126 Мбит/с и отдаче 108 Мбит/с видеотрафика в сторону коммутатора ядра. 
    \item \textbf{Сервер автоматизированной сборки (CI/CD):} Для компиляции тестового программного обеспечения и его одновременной загрузки на тестируемые устройства выбран сервер Гравитон С2122АА. Вычислительная платформа оснащена сетевыми адаптерами 10G SFP+, что гарантирует стабильную передачу пикового трафика прошивок (до 3900 Мбит/с для 78 IoT-устройств) без потерь и задержек. 
    \item \textbf{Сетевое хранилище (NAS):} Для кэширования дистрибутивов, хранения компонентов сборки и ведения логов тестирования применяется система хранения данных Qtech серии QSRV. Накопитель поддерживает организацию аппаратного RAID-массива для повышения сохранности данных. Подключение к сети реализовано через оптический порт 10G SFP+ (или агрегированную группу портов) для повышения суммарной пропускной способности при обмене данными с сервером сборки и рабочими станциями.
\end{itemize}




\subsection{Расчет РоЕ-бюджета}

Расчет энергетического баланса производится для подтверждения того, что источники питания выбранных коммутаторов способны одновременно запитать все подключенные по витой паре устройства на максимальной мощности с сохранением нормативного технологического запаса. В соответствии с правилами проектирования СКС вводится коэффициент запаса мощности 30\,\% ($K_{\text{зап}} = 1{,}3$), применяемый для компенсации омических потерь в протяженных трассах и учета пусковых токов. 

В связи с существенным различием потребляемой мощности оконечного оборудования, дистанционное питание в кроссовой разделено на две независимые подсистемы на базе разных коммутаторов.

\textbf{Подсистема питания видеокамер и СКУД (коммутатор Eltex MES2448P, PoE+)}

Суммарная потребляемая мощность оконечного оборудования рассчитывается исходя из максимальных паспортных данных устройств:

\begin{itemize}
    \item Узконаправленные IP-камеры QVC-MiR501Z1: $6 \times 8 \text{ Вт} = 48 \text{ Вт}$.
    \item Обзорные IP-камеры QVC-MiR803M: $2 \times 7 \text{ Вт} = 14 \text{ Вт}$.
    \item Контроллеры СКУД Sigur E4 PRO: $2 \times 4{,}5 \text{ Вт} = 9 \text{ Вт}$.
\end{itemize}
Базовая мощность подсистемы составляет:

$P_{\text{PoE+}} = 48 + 14 + 9 = 71 \text{ Вт}$.

С учетом нормативного коэффициента запаса 30\,\% требуемый бюджет мощности для данной группы равен: $P_{\text{PoE+\_треб}} = 71 \times 1{,}3 = 92{,}3 \text{ Вт}$. 

Аппаратный PoE-бюджет выбранного коммутатора Eltex MES2448P составляет 720 Вт. Это значение с многократным избытком перекрывает расчетную потребность в 92,3 Вт, гарантируя стабильную подачу питания на камеры и СКУД.

\textbf{Подсистема питания моноблоков операторов (коммутатор Aquarius, PoE++)}

Питание 6 моноблоков «Гравитон» М75И осуществляется от высокомощных портов коммутатора Aquarius AQ-N3000-48P4Y2Q. Расчет ведется исходя из максимального заявленного энергопотребления моноблока (120 Вт).

Базовая нагрузка постов операторской составляет:

$P_{\text{PoE++}} = 6 \times 120 \text{ Вт} = 720 \text{ Вт}$.

С учетом технологического запаса 30\,\% требуемый бюджет мощности для безотказной работы постов равен: $P_{\text{PoE++\_треб}} = 720 \times 1{,}3 = 936 \text{ Вт}$. 

Номинальный базовый PoE-бюджет коммутатора Aquarius составляет 1080 Вт. Устройство полностью покрывает расчетную нагрузку в 936 Вт. (Примечание: поскольку пиковая потребляемая мощность моноблока в 120 Вт превышает лимиты стандарта IEEE 802.3bt Type 4, ограниченные значением 90 Вт на порт, для корректной работы моноблоков порты коммутатора Aquarius должны аппаратно поддерживать и быть сконфигурированы на выдачу проприетарной мощности Ultra PoE до 120 Вт, либо рабочее потребление самих моноблоков должно программно ограничиваться профилем энергосбережения операционной системы).

\section{Выбор источников бесперебойного питания}

Для обеспечения автономной работы активного сетевого и серверного оборудования в течение 30 минут при отключении внешнего электроснабжения проектом предусмотрена установка источников бесперебойного питания (ИБП) с топологией On-Line (двойное преобразование). 

Расчет требуемой мощности ИБП базируется на суммарном пиковом энергопотреблении устройств. Для компенсации пусковых токов и обеспечения эксплуатационной надежности вводится нормативный коэффициент запаса $K_{\text{зап}} = 1{,}2$.

\textbf{Расчет для телекоммуникационного шкафа серверной (пом. №\,19)}

Пиковое энергопотребление вычислительного оборудования серверной комнаты составляет:
\begin{itemize}
    \item Медиасервер (Гравитон С2242И): 1600 Вт;
    \item Сервер автоматизированной сборки (Гравитон С2122АА): 1600 Вт;
    \item Сетевое хранилище (Qtech QSRV): 850 Вт;
    \item Сервисный маршрутизатор (Eltex ESR-1200): 85 Вт;
    \item Коммутатор ядра (Eltex MES5324): 107 Вт.
\end{itemize}

Суммарная активная потребляемая мощность серверного шкафа составляет:
\begin{equation}
    P_{\text{серв}} = 1600 + 1600 + 850 + 85 + 107 = 4242 \text{ Вт}.
\end{equation}

Требуемая активная мощность ИБП с учетом коэффициента запаса:
\begin{equation}
    P_{\text{ИБП\_серв}} = 4242 \times 1{,}2 = 5090{,}4 \text{ Вт}.
\end{equation}

Для защиты оборудования серверной выбран стоечный ИБП Сайбер Электро РСК-Эксперт-II-6000С-М (мощность 6000 ВА / 6000 Вт, коэффициент мощности $\cos\varphi = 1{,}0$). Данная модель полностью покрывает расчетную нагрузку и обладает усиленным зарядным устройством (до 8 А) для работы с внешними массивами аккумуляторов.

\textbf{Расчет для телекоммуникационного шкафа кроссовой (пом. №\,16)}

Пиковое энергопотребление стека коммутаторов доступа (с учетом полной нагрузки PoE) составляет:
\begin{itemize}
    \item Коммутатор подсистемы PoE++ (Aquarius AQ-N3000-48P4Y2Q): 2200 Вт;
    \item Коммутатор подсистемы PoE+ (Eltex MES2448P): 820 Вт;
    \item Коммутаторы базовые без PoE (Eltex MES2448B, 2 шт.): $2 \times 66 \text{ Вт} = 132 \text{ Вт}$.
\end{itemize}

Суммарная активная потребляемая мощность шкафа кроссовой:
\begin{equation}
    P_{\text{кросс}} = 2200 + 820 + 132 = 3152 \text{ Вт}.
\end{equation}

Требуемая активная мощность ИБП для кроссовой комнаты:
\begin{equation}
    P_{\text{ИБП\_кросс}} = 3152 \times 1{,}2 = 3782{,}4 \text{ Вт}.
\end{equation}

Для унификации инфраструктуры в кроссовой комнате применяется аналогичный ИБП Сайбер Электро РСК-Эксперт-II-6000. Для корректного завершения работы серверов (Graceful Shutdown) оба ИБП оснащаются платами сетевого мониторинга СЭ-SNMP-202.

\section{Расчет емкости аккумуляторных батарей (АКБ)}

Время автономной работы системы должно составлять не менее 30 минут ($T = 0{,}5$ ч). Требуемая полезная емкость аккумуляторных батарей рассчитывается с учетом КПД инвертора ИБП ($\eta_{\text{инв}} = 0{,}95$) и рабочего напряжения шины постоянного тока ($U_{\text{шина}} = 240$ В).

Расчет минимальной требуемой емкости для серверного шкафа (при максимальной нагрузке 4242 Вт):
\begin{equation}
    C_{\text{акб\_мин}} = \frac{P_{\text{серв}} \cdot T}{\eta_{\text{инв}} \cdot U_{\text{шина}}} = \frac{4242 \cdot 0{,}5}{0{,}95 \cdot 240} = \frac{2121}{228} \approx 9{,}3 \text{ А}\cdot\text{ч}.
\end{equation}

Для обеспечения заданного времени работы с учетом старения свинцово-кислотных элементов (деградация емкости до 40\,\% к концу срока службы) и ограничения глубины разряда, для ИБП серверной выбран напольный батарейный кабинет БШ-20, укомплектованный 20 промышленными АКБ емкостью 55 А$\cdot$ч. Меньшая нагрузка шкафа кроссовой (3152 Вт) позволяет обеспечить 30 минут автономии установкой одного стандартного стоечного батарейного модуля высотой 3U.

\section{Выбор системы климат-контроля и вентиляции}

В соответствии с требованиями ГОСТ Р 53246-2008, в телекоммуникационных помещениях поддерживается микроклимат в диапазоне температур от 18~$^\circ$C до 24~$^\circ$C. Мощность системы охлаждения рассчитывается на основе суммарного выделения тепла оборудованием с коэффициентом запаса 30\,\% ($K = 1{,}3$). Тепловые потери ИБП принимаются равными 5\,\% от полезной нагрузки.

\textbf{Охлаждение серверной комнаты (пом. №\,19)}

Тепловыделение серверного оборудования эквивалентно его энергопотреблению. Базовое тепловыделение шкафа составляет:
\begin{equation}
    Q_{\text{серв}} = P_{\text{серв}} + (P_{\text{серв}} \times 0{,}05) = 4242 + 212{,}1 = 4454{,}1 \text{ Вт}.
\end{equation}

Требуемая холодопроизводительность кондиционеров:
\begin{equation}
    Q_{\text{конд\_треб}} = 4454{,}1 \times 1{,}3 = 5790{,}3 \text{ Вт}.
\end{equation}

Для отвода тепла выбраны два внутристоечных прецизионных кондиционера C3 Solutions C3.AIR.BIS007 с номинальной холодопроизводительностью 7{,}2 кВт каждый. Система работает по схеме 100\,\% резервирования (N+1) с автоматической ротацией блоков. Электропитание кондиционеров осуществляется от силового щита здания независимыми линиями, в обход ИБП шкафа.

\textbf{Вентиляция кроссовой комнаты (пом. №\,16)}

Фактическая мощность, отдаваемая внешним PoE-потребителям (моноблокам и камерам), составляет 791 Вт. Данная энергия рассеивается вне телекоммуникационного шкафа. 
Внутреннее тепловыделение оборудования кроссовой равно:
\begin{equation}
    Q_{\text{кросс}} = (3152 - 791) + (3152 \times 0{,}05) = 2361 + 157{,}6 = 2518{,}6 \text{ Вт}.
\end{equation}

Для поддержания допустимого температурного режима без установки прецизионного кондиционера шкаф кроссовой комнаты оборудуется перфорированными дверьми и активным вентиляторным модулем на 4 элемента в крыше (работающим на вытяжку). Свободные юниты закрываются глухими фальш-панелями для исключения рециркуляции горячих потоков воздуха.


\section{Выбор моделей телекоммуникационных шкафов}

Для централизованного размещения активного сетевого оборудования, пассивных коммутационных элементов и систем обеспечения (ИБП, климат-контроль) в серверной (пом. №\,19) и кроссовой (пом. №\,16) комнатах проектом предусмотрена установка напольных 19-дюймовых телекоммуникационных шкафов. Выбор моделей базируется на габаритах активного оборудования, требованиях к отводу тепла и допустимой нагрузке на железобетонные перекрытия.

Конструкции выбранных шкафов отвечают следующим общим эксплуатационным требованиям:
\begin{itemize}
    \item Наличие съемных боковых панелей с индивидуальными замками для оперативного доступа к вертикальным кабельным органайзерам и облегчения монтажа.
    \item Наличие перфорации на дверях для обеспечения беспрепятственного сквозного воздушного потока (Front-to-Back) и пассивного отвода тепла.
    \item Обеспечение контура заземления: все металлические элементы каркаса, панели и двери объединяются проводниками системы уравнивания потенциалов и подключаются к главной заземляющей шине здания в соответствии с ГОСТ Р 50571.5.54.
\end{itemize}

\textbf{Расчет массогабаритных параметров и нагрузки на перекрытия}

Важным критерием проектирования является соблюдение предельно допустимой равномерно распределенной нагрузки на существующие межэтажные перекрытия, которая по строительным нормам составляет 400 кг/м$^2$. 
Максимально допустимый вес внутреннего оборудования ($M_{\text{max}}$) для конкретного шкафа вычисляется по формуле:
\begin{equation}
    M_{\text{max}} = 400 \cdot W \cdot L - M_{\text{шкаф}},
\end{equation}
где:
\begin{itemize}
    \item $W$ --- внешняя ширина шкафа, м;
    \item $L$ --- внешняя глубина шкафа, м;
    \item $M_{\text{шкаф}}$ --- собственный вес пустого шкафа, кг.
\end{itemize}

\subsection{Оборудование серверной комнаты (пом. №\,19)}

Для размещения вычислительной группы (медиасервер, сервер CI/CD, СХД), маршрутизатора ядра, двух внутристоечных кондиционеров и силового модуля ИБП выбран шкаф \textbf{C3 Solutions серии C3.RK (артикул: C3.RK426010.2P2P)}.
\begin{itemize}
    \item \textbf{Габариты и емкость:} Высота 42U (2010 мм), ширина 600 мм, внешняя глубина 1000 мм (полезная монтажная глубина 900 мм, что оптимально для серверов).
    \item \textbf{Конструкция дверей:} Двустворчатые стальные двери с развитой перфорацией 80\,\%, что превосходит базовые нормативы (60--70\,\%) и гарантирует отсутствие тепловых мешков при работе мощных вычислительных узлов.
    \item \textbf{Грузоподъемность и масса:} Статическая несущая способность усиленного каркаса составляет 1500 кг, собственный вес пустого шкафа --- 138 кг.
\end{itemize}

\textit{Расчет нагрузки на перекрытие для серверного шкафа:}

Базовая площадь опоры шкафа составляет $S = 0{,}6 \times 1{,}0 = 0{,}6 \text{ м}^2$.
Допустимая общая масса шкафа с оборудованием: $400 \times 0{,}6 = 240 \text{ кг}$.
Лимит массы для активного оборудования: $M_{\text{max\_серв}} = 240 - 138 = 102 \text{ кг}$.
Фактическая масса планируемого к установке оборудования (три сервера, коммутатор, маршрутизатор, два кондиционера и силовой блок ИБП без батарей) составляет около 155 кг. Так как расчетная масса превышает допустимый лимит ($155 \text{ кг} > 102 \text{ кг}$), для соблюдения строительных нормативов проектом предусмотрена установка данного шкафа на металлическую распределительную опорную раму площадью $0{,}8 \text{ м}^2$. Это увеличивает общий лимит массы до 320 кг и позволяет безопасно разместить оборудование весом до 182 кг, удовлетворяя требования нормативных документов.

\subsection{Оборудование кроссовой комнаты (пом. №\,16)}

Для размещения стека из четырех коммутаторов уровня доступа, патч-панелей, кабельных органайзеров и стоечного ИБП с аккумуляторным блоком выбран шкаф \textbf{C3 Solutions серии C3.RI (артикул: C3.RI3303)}.
\begin{itemize}
    \item \textbf{Габариты и емкость:} Высота 33U (1620 мм), ширина 600 мм, глубина 800 мм (монтажная глубина 700--720 мм). Данные габариты обеспечивают идеальный баланс между экономией пространства помещения и наличием места для укладки запасов витой пары.
    \item \textbf{Конструкция дверей:} Одинарные двери со стальными полотнами, обладающими высоким уровнем перфорации для сквозной пассивной вентиляции PoE-коммутаторов.
    \item \textbf{Грузоподъемность и масса:} Несущая способность составляет 800 кг, собственный вес пустого шкафа --- 90 кг.
\end{itemize}

\textit{Расчет нагрузки на перекрытие для кроссового шкафа:}

Площадь опоры шкафа: $S = 0{,}6 \times 0{,}8 = 0{,}48 \text{ м}^2$.
Допустимая общая масса шкафа с оборудованием: $400 \times 0{,}48 = 192 \text{ кг}$.
Лимит массы для размещаемого оборудования: $M_{\text{max\_кросс}} = 192 - 90 = 102 \text{ кг}$.
Фактическая масса оборудования кроссовой (4 коммутатора, ИБП с батарейным блоком 3U, панели) не превышает 70 кг. Условие $70 \text{ кг} \le 102 \text{ кг}$ строго выполняется, шкаф устанавливается непосредственно на пол без применения распределительных конструкций.

\section{Схемы компоновки телекоммуникационных шкафов}

Схема компоновки телекоммуникационных шкафов строится в строгом соответствии с принципом: тяжелое оборудование размещается внизу, оборудование с высоким тепловыделением --- вверху, а между коммутаторами и патч-панелями обязательно монтируются кабельные организаторы.

\subsection{Схема компоновки серверного шкафа (Помещение №\,19)}

\begin{itemize}
    \item \textbf{Верхняя зона:} На крыше шкафа устанавливается вентиляторный модуль, работающий на выдув горячего воздуха. В самых верхних юнитах размещается наиболее легкое оборудование: главный маршрутизатор, а под ним монтируется блок электрических розеток (БРП) для питания верхнего эшелона устройств.
    \item \textbf{Зона коммутации ядра:} Ниже устанавливается оптическая кросс-панель для терминирования внешних оптических линий, горизонтальный кабельный органайзер и коммутатор ядра сети. Жгуты кабелей разводятся по боковым вертикальным органайзерам, чтобы не перекрывать вентиляцию задней стенки.
    \item \textbf{Нижняя зона (вычислительная):} В нижней части шкафа размещается тяжелое серверное оборудование. Устанавливаются сервер CI/CD, медиасервер и сетевое хранилище (NAS), которое монтируется ниже остальных серверов из-за большого веса массива накопителей.
    \item \textbf{Основание шкафа:} В самом низу устанавливается массивный источник бесперебойного питания (ИБП) и нижний блок электрических розеток.
    \item \textbf{Резервное пространство:} Все неиспользуемые посадочные места (юниты) на лицевой стороне закрываются глухими фальш-панелями, что блокирует переток горячего воздуха от задней стенки к передним воздухозаборникам активного оборудования.
\end{itemize}

\begin{figure}[H]
    \centering
    \resizebox{!}{0.85\textheight}{
    \begin{tikzpicture}[x=1.2cm, y=0.5cm, font=\sffamily\scriptsize] 
        \tikzset{
            net/.style={draw=black, fill=red!20, minimum width=8.4cm, minimum height=0.5cm, align=center},
            cool/.style={draw=black, fill=blue!20, minimum width=8.4cm, minimum height=0.5cm, align=center},
            srv/.style={draw=black, fill=teal!20, minimum width=8.4cm, minimum height=0.5cm, align=center},
            pwr/.style={draw=black, fill=yellow!30, minimum width=8.4cm, minimum height=0.5cm, align=center},
            pas/.style={draw=black, fill=gray!20, minimum width=8.4cm, minimum height=0.5cm, align=center},
            empty/.style={draw=black, fill=gray!5, minimum width=8.4cm, minimum height=0.5cm, align=center}
        }
        \def\drawunit#1#2#3#4{
            \pgfmathsetmacro{\ypos}{#1 - 1 + #2/2}
            \node[#3, minimum height=#2*0.5cm, inner sep=0pt] at (3.5, \ypos) {#4};
            \foreach \u in {#1,...,\the\numexpr#1+#2-1\relax} {
                \node[anchor=east, font=\scriptsize, text=gray] at (-0.2, \u-0.5) {\u};
                \node[anchor=west, font=\scriptsize, text=gray] at (7.2, \u-0.5) {\u};
            }
        }
        \draw[line width=1pt] (0,0) rectangle (7, 42);
        
        \drawunit{1}{1}{empty}{Резерв}
        \drawunit{2}{2}{pwr}{\textbf{Нижний БРП (PDU)}\\(2U)}
        \drawunit{4}{4}{pwr}{\textbf{ИБП Сайбер Электро РСК-Эксперт-II-6000С-М}\\(6000 Вт, 4U)}
        \drawunit{8}{2}{srv}{\textbf{СХД NAS Qtech QSRV}\\(850 Вт, 2U)}
        \drawunit{10}{2}{srv}{\textbf{Сервер CI/CD Гравитон С2122АА}\\(1600 Вт, 2U)}
        \drawunit{12}{2}{srv}{\textbf{Медиасервер Гравитон С2242И}\\(1600 Вт, 2U)}
        
        \drawunit{14}{14}{empty}{Глухие фальш-панели (Свободный резерв 14U)}
        
        \drawunit{28}{5}{cool}{\textbf{Кондиционер C3.AIR.BIS007 (Резервный)}\\(Холодопроизв. 7{,}2 кВт, 5U)}
        \drawunit{33}{5}{cool}{\textbf{Кондиционер C3.AIR.BIS007 (Основной)}\\(Холодопроизв. 7{,}2 кВт, 5U)}
        
        \drawunit{38}{1}{net}{\textbf{Коммутатор ядра Eltex MES5324} (107 Вт, 1U)}
        \drawunit{39}{1}{pas}{Горизонтальный кабельный органайзер (1U)}
        \drawunit{40}{1}{pas}{Оптическая кросс-панель (1U)}
        \drawunit{41}{1}{net}{\textbf{Маршрутизатор Eltex ESR-1200} (85 Вт, 1U)}
        \drawunit{42}{1}{pwr}{\textbf{Верхний БРП (PDU)} (1U)}
    
        \node[anchor=north] at (3.5, -0.8) {
            \begin{tabular}{ccccc}
            \tikz\node[net, minimum width=0.8cm, minimum height=0.4cm]{}; Сеть &
            \tikz\node[srv, minimum width=0.8cm, minimum height=0.4cm]{}; Серверы/СХД &
            \tikz\node[cool, minimum width=0.8cm, minimum height=0.4cm]{}; Охлаждение &
            \tikz\node[pwr, minimum width=0.8cm, minimum height=0.4cm]{}; Питание &
            \tikz\node[pas, minimum width=0.8cm, minimum height=0.4cm]{}; Пассивное
            \end{tabular}
        };
    \end{tikzpicture}
    }
    \caption{Схема компоновки серверного шкафа (42U)}
\end{figure}

\subsection{Размещение оборудования в кроссовой комнате}
Для размещения оборудования этажного распределительного пункта (ЭРП) в кроссовой комнате (пом. №\,16) выбран напольный телекоммуникационный шкаф C3 Solutions серии C3.RI высотой 33U. В отличие от серверной, оборудование уровня доступа характеризуется меньшим тепловыделением и не требует установки прецизионных внутристоечных кондиционеров. Охлаждение осуществляется за счет сквозной перфорации дверей и активного вентиляторного модуля, установленного в крыше шкафа (работает на вытяжку горячего воздуха).

Компоновка оборудования по вертикали (сверху вниз) выполнена с учетом распределения весовых нагрузок, тепловых потоков и высокой плотности кабельных соединений (High Density):

\begin{itemize}
    \item \textbf{Верхняя зона:} В самых верхних юнитах шкафа, непосредственно под вентиляторным модулем, монтируется блок розеток (БРП) для питания верхнего эшелона устройств. Ниже устанавливается оптическая кросс-панель высокой плотности для терминирования волоконно-оптических магистралей (10G SFP+), приходящих от ядра сети из серверной комнаты.
    \item \textbf{Зона коммутации (уровень доступа):} Центральная часть шкафа отведена под терминацию 150 абонентских кабельных линий (витая пара Cat.6A). Поскольку к оборудованию подводится массивный жгут кабелей, применено строгое правило чередования: между каждой 24-портовой медной патч-панелью и каждым коммутатором устанавливается горизонтальный кабельный органайзер высотой 1U.
    В этой зоне последовательно размещаются:
    \begin{itemize}
        \item Коммутатор подсистемы PoE++ (Aquarius AQ-N3000-48P4Y2Q) для питания моноблоков операторов;
        \item Коммутатор подсистемы PoE+ (Eltex MES2420-48P) для обслуживания IP-камер и контроллеров СКУД;
        \item Два базовых коммутатора (Eltex MES2448B) для подключения стендов с IoT-устройствами и приборов КИП.
    \end{itemize}
    Входящие из-за стен жгуты кабелей разводятся исключительно по вертикальным боковым органайзерам шкафа. Это решение предотвращает перекрытие задних воздуховодов оборудования, исключает изломы кабелей и минимизирует механическую нагрузку на порты RJ-45.
    \item \textbf{Свободная зона (резерв):} Все неиспользуемые посадочные места между устройствами на лицевой стороне шкафа закрываются глухими металлическими фальш-панелями. Данная мера необходима для разделения холодных и горячих коридоров: панели физически блокируют переток отработанного горячего воздуха от задней стенки к фронтальным воздухозаборникам коммутаторов.
    \item \textbf{Нижняя зона:} В соответствии с правилом снижения центра тяжести («тяжелое --- вниз»), в самом основании шкафа (нижние 6U) устанавливается силовое оборудование: источник бесперебойного питания (ИБП) и внешний батарейный модуль (АКБ), обеспечивающие автономность работы подсистемы PoE и сетевого доступа. На уровне 1U закреплен нижний БРП.
\end{itemize}

В целях обеспечения электробезопасности и электромагнитной совместимости, все металлические элементы конструкции шкафа (рама, двери, боковые панели), а также корпуса коммутаторов и экранированные патч-панели соединяются проводниками защитного заземления (желто-зеленого цвета) и подключаются к главной заземляющей шине здания. При этом экраны абонентских кабелей заземляются исключительно на стороне патч-панелей кроссового шкафа, что исключает образование паразитных «земляных петель».

\begin{figure}[H]
    \centering
    \resizebox{!}{0.7\textheight}{
    \begin{tikzpicture}[x=1.2cm, y=0.5cm, font=\sffamily\scriptsize]
        \tikzset{
            net/.style={draw=black, fill=red!20, minimum width=8.4cm, minimum height=0.5cm, align=center},
            cool/.style={draw=black, fill=blue!20, minimum width=8.4cm, minimum height=0.5cm, align=center},
            srv/.style={draw=black, fill=teal!20, minimum width=8.4cm, minimum height=0.5cm, align=center},
            pwr/.style={draw=black, fill=yellow!30, minimum width=8.4cm, minimum height=0.5cm, align=center},
            pas/.style={draw=black, fill=gray!20, minimum width=8.4cm, minimum height=0.5cm, align=center},
            empty/.style={draw=black, fill=gray!5, minimum width=8.4cm, minimum height=0.5cm, align=center}
        }
        \def\drawunit#1#2#3#4{
            \pgfmathsetmacro{\ypos}{#1 - 1 + #2/2}
            \node[#3, minimum height=#2*0.5cm, inner sep=0pt] at (3.5, \ypos) {#4};
            \foreach \u in {#1,...,\the\numexpr#1+#2-1\relax} {
                \node[anchor=east, font=\scriptsize, text=gray] at (-0.2, \u-0.5) {\u};
                \node[anchor=west, font=\scriptsize, text=gray] at (7.2, \u-0.5) {\u};
            }
        }
        \draw[line width=1pt] (0,0) rectangle (7, 33);
        
        \drawunit{1}{1}{empty}{Резерв}
        \drawunit{2}{3}{pwr}{\textbf{Доп. батарейный модуль ИБП (АКБ)}\\(3U)}
        \drawunit{5}{4}{pwr}{\textbf{ИБП Сайбер Электро РСК-Эксперт-II-6000}\\(4U)}
        
        \drawunit{9}{4}{empty}{Глухие фальш-панели (Свободный резерв 4U)}
        
        \drawunit{13}{1}{pas}{Горизонтальный кабельный органайзер (1U)}
        \drawunit{14}{1}{net}{\textbf{Коммутатор Eltex MES2448B} (1U)}
        \drawunit{15}{1}{pas}{Горизонтальный кабельный органайзер (1U)}
        \drawunit{16}{1}{pas}{Медная патч-панель 24 порта (1U)}
        
        \drawunit{17}{1}{pas}{Горизонтальный кабельный органайзер (1U)}
        \drawunit{18}{1}{net}{\textbf{Коммутатор Eltex MES2448B} (1U)}
        \drawunit{19}{1}{pas}{Горизонтальный кабельный органайзер (1U)}
        \drawunit{20}{1}{pas}{Медная патч-панель 24 порта (1U)}
        
        \drawunit{21}{1}{pas}{Горизонтальный кабельный органайзер (1U)}
        \drawunit{22}{1}{net}{\textbf{Коммутатор PoE+ Eltex MES2420-48P} (1U)}
        \drawunit{23}{1}{pas}{Горизонтальный кабельный органайзер (1U)}
        \drawunit{24}{1}{pas}{Медная патч-панель 24 порта (1U)}
        
        \drawunit{25}{1}{pas}{Горизонтальный кабельный органайзер (1U)}
        \drawunit{26}{1}{net}{\textbf{Коммутатор PoE++ Aquarius} (1U)}
        \drawunit{27}{1}{pas}{Горизонтальный кабельный органайзер (1U)}
        \drawunit{28}{1}{pas}{Медная патч-панель 24 порта (1U)}
        
        \drawunit{29}{2}{empty}{Глухие фальш-панели (Резерв 2U)}
        
        \drawunit{31}{1}{pas}{Оптическая патч-панель (1U)}
        \drawunit{32}{1}{pwr}{\textbf{Верхний распределительный щит / БРП} (1U)}
        \drawunit{33}{1}{empty}{Резерв (1U)}
    
        \node[anchor=north] at (3.5, -0.8) {
            \begin{tabular}{cccc}
            \tikz\node[net, minimum width=0.8cm, minimum height=0.4cm]{}; Сеть &
            \tikz\node[pwr, minimum width=0.8cm, minimum height=0.4cm]{}; Питание &
            \tikz\node[pas, minimum width=0.8cm, minimum height=0.4cm]{}; Пассивное
            \end{tabular}
        };
    \end{tikzpicture}
    }
    \caption{Схема компоновки кроссового шкафа (33U)}
\end{figure}

% ======================================================================
%  ОСНОВНАЯ ЧАСТЬ
% ======================================================================

\chapter{Обследование объекта}

\section{Общие сведения об объекте}
Текст подраздела\ldots

\section{Разработка архитектурного плана объекта}
Текст подраздела\ldots

\section{Размещение объектов инфраструктуры}
Текст подраздела\ldots

\section{Подсчет количества конечных устройств}
Текст подраздела\ldots


\chapter{Необходимое сетевое оборудование}

\section{Расчет пропускной способности каналов связи}
Текст подраздела\ldots

\section{Архитектура сети}
Текст подраздела\ldots

\section{Обоснование выбора коммутаторов и сетевого оборудования}
Текст подраздела\ldots


\chapter{Устройство и выбор шкафа}

\section{Выбор патч-панели}
Текст подраздела\ldots

\section{Выбор органайзера}
Текст подраздела\ldots

\section{Выбор БРП}
Текст подраздела\ldots

\section{Обоснование выбора ИБП}
Текст подраздела\ldots

\section{Обоснование выбора охлаждающего элемента и систем вентиляции}
Текст подраздела\ldots

\section{Выбор АКБ}
Текст подраздела\ldots

\section{Обоснование выбора шкафа}
Текст подраздела\ldots

\section{Расчет масс и габаритов шкафов}
Текст подраздела\ldots

\section{Расчет заполнения шкафов}
Текст подраздела\ldots

\section{Выбор места размещения шкафа}
Текст подраздела\ldots


\chapter{Выбор и расчет кабеля}

\section{Обоснование выбора типов кабелей}
Текст подраздела\ldots

\section{Расчет длины кабеля}
Текст подраздела\ldots

\section{Расчет длины кабеля от перфорированного лотка до шкафа}
Текст подраздела\ldots

\section{Расчет максимальной длины кабеля и проверка выбора}
Текст подраздела\ldots

\section{Расчет закупочного количества бухт кабеля}
Текст подраздела\ldots

\section{Сводная таблица результатов расчета требуемой длины кабелей}
Текст подраздела\ldots


\chapter{Расчет элементов кабельной инфраструктуры}

\section{Выбор и расчет количества розеток}
Текст подраздела\ldots

\section{Расчет диаметра и площади сечения пучков кабелей}
Текст подраздела\ldots

\section{Расчет геометрических размеров и выбор лотков}
Текст подраздела\ldots

\section{Расчет геометрических размеров и выбор кабель-каналов}
Текст подраздела\ldots

\section{Выбор элементов крепления}
Текст подраздела\ldots

\section{Расчет муфт, используемых для проводки кабелей через стены}
Текст подраздела\ldots

\section{Заземление шкафов, лотков и экранированных кабелей}
Текст подраздела\ldots


\chapter{Расчет оптоволоконной линии}

\section{Выбор оптоволоконного кабеля, компонентов соединений}
Текст подраздела\ldots

\section{Расчет длины кабеля (9, 11 этажи)}
Текст подраздела\ldots

\section{Прокладка оптоволоконного кабеля}
Текст подраздела\ldots


\chapter{Описание маркировки проводов}
Текст раздела\ldots


% ======================================================================
%  ЗАКЛЮЧЕНИЕ
% ======================================================================
\chapter*{ЗАКЛЮЧЕНИЕ}
\addcontentsline{toc}{chapter}{ЗАКЛЮЧЕНИЕ}

Текст заключения\ldots

% ======================================================================
%  СПИСОК ИСПОЛЬЗОВАННЫХ ИСТОЧНИКОВ
% ======================================================================
%\makebibliography

% ======================================================================
%  ПРИЛОЖЕНИЕ А --- ТЕХНИЧЕСКОЕ ЗАДАНИЕ
% ======================================================================
\begin{appendices}
\addtocontents{toc}{\protect\setcounter{tocdepth}{0}}
\setcounter{section}{0}
\renewcommand{\thesection}{\arabic{section}}

\cleardoublepage
\stepcounter{chapter}
\phantomsection
\addcontentsline{toc}{chapter}{ПРИЛОЖЕНИЕ А}

% --- Титульная страница ТЗ ---
\thispagestyle{empty}
\begin{flushright}
ПРИЛОЖЕНИЕ А
\end{flushright}

\vspace*{0pt}
\noindent
\begin{minipage}[t]{0.48\textwidth}
    \raggedright
    \textbf{СОГЛАСОВАНО}\\[6pt]
    Руководитель НИР\\[10pt]
    \blank[3cm]~И.М. Сидякин\\[6pt]
    <<\blank[0.6cm]>>~\blank[2.5cm]~2026~г.
\end{minipage}%
\hfill
\begin{minipage}[t]{0.48\textwidth}
    \raggedright
    \textbf{УТВЕРЖДАЮ}\\[6pt]
    Заведующий кафедрой ИУ3\\[10pt]
    \blank[3cm]~А.Н. Алфимцев\\[6pt]
    <<\blank[0.6cm]>>~\blank[2.5cm]~2026~г.
\end{minipage}

\vfill
\begin{center}
    {\LARGE \bfseries ТЕХНИЧЕСКОЕ ЗАДАНИЕ НА НИР}\\[8mm]
    {\Large \bfseries
    ПРОЕКТИРОВАНИЕ СТРУКТУРИРОВАННОЙ КАБЕЛЬНОЙ\\
    СИСТЕМЫ ОФИСА ЛАБОРАТОРИИ\\
    ТЕСТИРОВАНИЯ УСТРОЙСТВ IoT}
\end{center}
\vfill

\noindent
\textbf{РАЗРАБОТАНО}\\[6pt]
Студент ИУ3-33Б\\[10pt]
\blank[3cm]~А.А. Захватов\\[6pt]
<<\blank[0.6cm]>>~\blank[2.5cm]~2026~г.

\vspace{20mm}

\begin{center}
    Москва\\
    2026~г.
\end{center}

\newpage

% =======================================================================
%  РАЗДЕЛ 1. ОБЩИЕ СВЕДЕНИЯ
% =======================================================================

\section{Полное наименование системы и ее условное обозначение}
\begin{enumerate}[label=\arabic{section}.\arabic*]
    \item \textbf{Полное наименование:} Структурированная кабельная система (СКС) этажного офиса лаборатории тестирования аппаратуры и программного обеспечения устройств IoT.
    \item \textbf{Условное обозначение:} СКС-IoT-ИУ3-V10.
\end{enumerate}

\section{Шифр темы или основание для проведения работ}
\begin{enumerate}[label=\arabic{section}.\arabic*]
    \item \textbf{Основание для разработки:} Индивидуальное задание к курсовой работе по дисциплине <<Проектирование структурированной кабельной системы>>, 3-й семестр бакалавриата кафедры ИУ-3.
\end{enumerate}

\section{Наименование предприятий (организаций) --- Заказчика и Разработчика}
\begin{enumerate}[label=\arabic{section}.\arabic*]
    \item \textbf{Заказчик:} Кафедра ИУ-3 МГТУ им. Н.Э. Баумана.
    \item \textbf{Разработчик:} Студент группы ИУ3-33Б, Захватов Антон Алексеевич.
\end{enumerate}

\section{Сроки начала и окончания работы}
\begin{enumerate}[label=\arabic{section}.\arabic*]
    \item \textbf{Начало:} Сентябрь 2026 г.
    \item \textbf{Окончание:} Декабрь 2026 г.
\end{enumerate}

% =======================================================================
%  РАЗДЕЛ 2. ЦЕЛИ И НАЗНАЧЕНИЕ СОЗДАНИЯ КАБЕЛЬНОЙ СИСТЕМЫ
% =======================================================================

\section{Назначение системы}
СКС предназначена для создания универсальной, надежной и масштабируемой кабельной инфраструктуры, обеспечивающей передачу данных, сигналов управления и электропитания по технологии PoE (Power over Ethernet) для обеспечения функционирования:
\begin{enumerate}[label=\asbuk*)]
    \item лабораторных стендов тестирования IoT-устройств и контрольно-измерительного оборудования;
    \item рабочих мест операторов и административного персонала;
    \item систем охранного видеонаблюдения контроля стендов и общих зон;
    \item систем контроля и управления доступом (СКУД);
    \item магистральных каналов связи с серверными здания.
\end{enumerate}

\section{Цели создания системы}
\begin{enumerate}[label=\arabic{section}.\arabic*]
    \item \textbf{Обеспечение требуемой пропускной способности} для передачи трафика от IoT-устройств и систем видеонаблюдения высокого разрешения (в режимах MCU и SFU).
    \item \textbf{Повышение надежности} за счет организации топологии <<иерархическая звезда>> и избыточности магистральных линий.
    \item \textbf{Оптимизация эксплуатации} ИТ-инфраструктуры за счет сведения всех кабельных систем к единому коммутационному центру.
    \item \textbf{Выполнение импортозамещения} путем подбора активного сетевого оборудования из Единого реестра российской радиоэлектронной продукции Минпромторга РФ.
\end{enumerate}

% =======================================================================
%  РАЗДЕЛ 3. ХАРАКТЕРИСТИКА КАБЕЛЬНОЙ СИСТЕМЫ
% =======================================================================

\section{Краткие сведения об объекте}
\begin{enumerate}[label=\arabic{section}.\arabic*]
    \item \textbf{Расположение:} Один этаж в многоэтажном здании.
    \item \textbf{Общая площадь:} 613,08~м$^2$ в осях.
    \item \textbf{Количество помещений:} 24 (согласно экспликации архитектурного плана формата А2).
    \item \textbf{Режим работы:} Непрерывный, с повышенной нагрузкой на кабельные тракты в зонах испытаний.
\end{enumerate}

\section{Перечень оборудуемых помещений и плотность рабочих мест}
Распределение информационных портов и оборудования выполняется строго по экспликации:

\begin{enumerate}[label=\arabic{section}.\arabic*]
    \item \textbf{Помещения лабораторий:}
    \begin{enumerate}[label=\arabic{section}.\arabic{enumi}.\arabic*]
        \item \textbf{Помещение №4 (Кабинет 3):} 3~лабораторных стенда. На каждом стенде: 18~портов IoT-устройств, 1~порт камеры контроля (PoE+), 6~резервных портов. Емкость MUTOA стенда = 25~портов.
        \item \textbf{Помещение №11 (Кабинет 1):} 3~лабораторных стенда. На каждом стенде: 8~портов IoT-устройств, 1~порт камеры контроля (PoE+), 6~резервных портов. Емкость MUTOA стенда = 15~портов.
        \item \textit{Дополнительно в №4 и №11:} по 1~камере общего вида (PoE+) и по 1~контроллеру СКУД (PoE+).
    \end{enumerate}

    \item \textbf{Операторская (Помещение №20 <<Опенспейс>>):}
    \begin{enumerate}[label=\arabic{section}.\arabic{enumi}.\arabic*]
        \item \textbf{Пост №1} (обслуживает 1~стенд): 1~моноблок (1~порт PoE++), 1~рабочая станция (1~порт), 1~резервный порт. Итого: 3~розетки.
        \item \textbf{Пост №2} (обслуживает 1~стенд): 1~моноблок (1~порт PoE++), 1~рабочая станция (1~порт), 1~резервный порт. Итого: 3~розетки.
        \item \textbf{Пост №3} (обслуживает 1~стенд): 1~моноблок (1~порт PoE++), 1~рабочая станция (1~порт), 1~резервный порт. Итого: 3~розетки.
        \item \textbf{Пост №4} (обслуживает 1~стенд): 1~моноблок (1~порт PoE++), 1~рабочая станция (1~порт), 1~резервный порт. Итого: 3~розетки.
        \item \textbf{Пост №5} (обслуживает 1~стенд): 1~моноблок (1~порт PoE++), 1~рабочая станция (1~порт), 1~резервный порт. Итого: 3~розетки.
        \item \textbf{Пост №6} (обслуживает 1~стенд): 1~моноблок (1~порт PoE++), 1~рабочая станция (1~порт), 1~резервный порт. Итого: 3~розетки.
        \item \textbf{Общая зона опенспейса:} 2~розетки СКС и 1~штатное место персонала (по 2~розетки).
    \end{enumerate}

    \item \textbf{Административные помещения:}
    \begin{enumerate}[label=\arabic{section}.\arabic{enumi}.\arabic*]
        \item Помещения №24 (Кабинет 2), №7 (Переговорная): оборудуются 2~розетками СКС каждое.
    \end{enumerate}

    \item \textbf{Технологические помещения:}
    \begin{enumerate}[label=\arabic{section}.\arabic{enumi}.\arabic*]
        \item \textbf{Помещение №19:} Серверная этажа (11,70~м$^2$). Размещение главного шкафа.
        \item \textbf{Помещение №16:} Кроссовая СКС (11,70~м$^2$). Размещение этажного распределительного пункта.
    \end{enumerate}
\end{enumerate}

% =======================================================================
%  РАЗДЕЛ 4. ТРЕБОВАНИЯ К КАБЕЛЬНОЙ СИСТЕМЕ
% =======================================================================

\section{Требования к структуре и архитектуре системы}
\begin{enumerate}[label=\arabic{section}.\arabic*]
    \item СКС должна быть спроектирована по архитектуре \textbf{иерархической звезды} с одним техническим центром (Серверная №19) и одной кроссовой комнатой (Помещение №16).
    \item В кроссовой СКС (№16) должен быть организован Этажный Распределительный Пункт (ЭРП).
    \item Магистральная подсистема должна связывать ЭРП этажа со слаботочным стояком (<<СКС>>) и организовывать вертикальные связи с Главной серверной здания (пом. 903, эт. 9) и Резервной серверной (пом. 1183, эт. 11).
\end{enumerate}

\section{Требования к производительности и компонентам}
\begin{enumerate}[label=\arabic{section}.\arabic*]
    \item Кабельная система горизонтальной подсистемы должна быть выполнена на основе медного четырехпарного кабеля <<витая пара>> категории не ниже \textbf{Cat.6A (класс Ea)} для поддержки скоростей до 10~Гбит/с и обеспечения стабильной работы PoE++/4PPoE (IEEE 802.3bt).
    \item Магистральная подсистема здания (вертикальная) должна быть спроектирована с использованием волоконно-оптического многомодового (OM4) или одномодового (OS2) кабеля для обеспечения отказоустойчивости и емкости каналов.
    \item Коммутационные панели (патч-панели) в шкафах должны иметь плотность не менее 24 портов RJ-45/1U с обязательной цветовой или цифровой маркировкой подсистем.
\end{enumerate}

\section{Требования к надежности и режимам видеонаблюдения}
Система должна обеспечивать бесперебойную передачу трафика от 8~камер (6~направленных, 2~обзорных) с поддержкой кодеков H.265/H.265+ в двух режимах работы медиасерверов:
\begin{enumerate}[label=\asbuk*)]
    \item штатный режим (MCU) --- архитектура с централизованным микшированием потоков;
    \item аварийный режим (SFU) --- архитектура с маршрутизацией сырых потоков напрямую на операторов, требующая проверки пиковых пропускных способностей каналов операторской подсистемы.
\end{enumerate}

\section{Требования к условиям эксплуатации и климатике}
\begin{enumerate}[label=\arabic{section}.\arabic*]
    \item В помещениях Серверной (№19) и Кроссовой (№16) должна быть предусмотрена принудительно-вытяжная вентиляция или кондиционирование для компенсации суммарного тепловыделения активного оборудования.
    \item Параметры микроклимата: температура 18--24~$^\circ$C, относительная влажность 30--55\% без конденсации.
    \item Active-оборудование шкафов должно питаться от выделенных источников бесперебойного питания (ИБП) с топологией On-Line и временем автономной работы при максимальной нагрузке не менее 30 минут.
\end{enumerate}

\section{Требования к кабельным трассам и конструкциям}
\begin{enumerate}[label=\arabic{section}.\arabic*]
    \item Магистральные лотки в коридорах должны быть металлическими (проволочными или перфорированными).
    \item Заполнение лотков и кабельных каналов не должно превышать \textbf{40--50}\% на этапе проектирования с учетом требований по радиусу изгиба и перспективного расширения.
    \item Расчет подвесов и элементов крепления должен проводиться с коэффициентом запаса прочности не менее 3 по отношению к весу заполненного лотка.
\end{enumerate}

\section{Требования к заземлению и электробезопасности}
\begin{enumerate}[label=\arabic{section}.\arabic*]
    \item Все металлические лотки, корпуса шкафов (19"), патч-панели и экраны кабелей должны быть интегрированы в единый контур защитного технологического заземления здания согласно ГОСТ Р 50571.5.54.
    \item Кабели должны отвечать требованиям пожарной безопасности по \textbf{ГОСТ 31565-2012} (исполнение не ниже нг-HF или нг-FRHF).
\end{enumerate}

% =======================================================================
%  РАЗДЕЛ 5. СОСТАВ И СОДЕРЖАНИЕ РАБОТ
% =======================================================================

\section{Состав и содержание работ по созданию кабельной системы}
В рамках курсового проектирования (согласно стадиям разработки по ГОСТ Р 15.301-2016) выполняются следующие этапы.
\begin{enumerate}[label=\arabic{section}.\arabic*]
    \item Анализ исходных данных. Изучение DXF-плана, расчет точного числа портов СКС для каждого помещения.
    \item Расчетно-теоретический этап:
    \begin{enumerate}[label=\asbuk*)]
        \item трассировка и расчет длин кабельных линий (метод контрольных длин);
        \item расчет сечения коробов, лотков, нагрузок на несущие конструкции;
        \item проектирование активных компонентов по реестру Минпромторга (коммутаторы доступа с PoE+/PoE++, агрегаторы, маршрутизаторы);
        \item расчет тепловыделения шкафов, параметров охлаждения и емкости аккумуляторных батарей ИБП.
    \end{enumerate}
    \item Графический этап. Создание схемы размещения рабочих мест на архитектурном плане (А2), составление схемы логической топологии и компоновки шкафов.
\end{enumerate}

% =======================================================================
%  РАЗДЕЛ 6. ПОРЯДОК КОНТРОЛЯ И ПРИЕМКИ
% =======================================================================

\section{Порядок контроля и приемки кабельной системы}

Контроль результатов учебной разработки осуществляет руководитель НИР.

При контроле проверяются:

\begin{itemize}
    \item соответствие проектного решения исходным данным индивидуального варианта;
    \item корректность размещения оборудования и кабельных трасс;
    \item корректность расчета количества портов, патч-панелей и заполнения шкафов;
    \item корректность расчетов трафика в режимах MCU и SFU;
    \item корректность расчетов кабельных линий, лотков, подвесов, охлаждения и ИБП;
    \item обоснованность выбора активного оборудования и PoE-бюджета;
    \item соответствие текстовой и графической документации требованиям задания и применимых нормативных документов.
\end{itemize}

В рамках учебной НИР фактический монтаж и натурные приемочные испытания не выполняются. При реализации проекта кабельные линии подлежат испытаниям и сертификационным измерениям в соответствии с требованиями выбранной кабельной системы и применимых стандартов.

% =======================================================================
%  РАЗДЕЛ 7. ТРЕБОВАНИЯ К ПОДГОТОВКЕ К ВВОДУ В ДЕЙСТВИЕ
% =======================================================================

\section{Требования к составу и содержанию работ по подготовке объекта к вводу кабельной системы в действие}

Перед вводом кабельной системы в эксплуатацию должны быть выполнены:

\begin{enumerate}[label=\asbuk*)]
    \item подготовка серверной №19 и кроссовой №16 к размещению телекоммуникационных шкафов;
    \item подготовка мест прокладки лотков, коробов и кабельных трасс;
    \item обеспечение электроснабжения, защитного заземления и охлаждения;
    \item монтаж и маркировка информационных розеток, MUTOA, патч-панелей и кабельных линий;
    \item установка и настройка активного оборудования;
    \item проверка кабельных линий и сетевых соединений;
    \item проверка работы PoE+/PoE++ оборудования и резервного питания;
    \item оформление исполнительной документации.
\end{enumerate}

% =======================================================================
%  РАЗДЕЛ 8. ТРЕБОВАНИЯ К ДОКУМЕНТИРОВАНИЮ
% =======================================================================

\section{Требования к документированию}

По результатам НИР должны быть подготовлены:

\begin{enumerate}[label=\asbuk*)]
    \item расчётно-пояснительная записка объёмом от 60 до 100 листов;
    \item техническое задание на кабельную систему;
    \item чертёж схемы сети объекта на листе формата A2;
    \item схема логической топологии сети;
    \item ведомость покупных изделий;
    \item расчёты портов, патч-панелей, шкафов, трафика, кабельных трасс, охлаждения и ИБП.
\end{enumerate}

% =======================================================================
%  РАЗДЕЛ 9. ИСТОЧНИКИ РАЗРАБОТКИ
% =======================================================================

\section{Источники разработки}

При разработке используются:

\begin{enumerate}[label=\asbuk*)]
    \item индивидуальное задание на выполнение НИР по теме <<Проектирование структурированной кабельной системы офиса лаборатории тестирования устройств IoT>>;
    \item исходные данные и архитектурный план индивидуального варианта;
    \item DXF-подоснова индивидуального варианта;
    \item ГОСТ Р 70443-2022, устанавливающий требования к техническому заданию на кабельную систему;
    \item ГОСТ 34.602-2020 и ГОСТ Р 15.301-2016, указанные в индивидуальном задании;
    \item ГОСТ 7.32-2017 в части оформления расчетно-пояснительной записки;
    \item ГОСТ Р 53246-2025 и ГОСТ Р 53245-2008 в части проектирования, монтажа и испытаний структурированных кабельных систем;
    \item ГОСТ Р 50571.5.54 в части защитного заземления;
    \item ГОСТ 31565-2012 в части пожарной безопасности кабельных изделий;
    \item Единый реестр российской радиоэлектронной продукции Минпромторга Российской Федерации.
\end{enumerate}

\cleardoublepage
\stepcounter{chapter}
\phantomsection
\addcontentsline{toc}{chapter}{ПРИЛОЖЕНИЕ Б}

\begin{flushright}
ПРИЛОЖЕНИЕ Б
\end{flushright}

\begin{center}
\textbf{\large Ведомость покупных изделий}
\end{center}

В данном разделе представлена сводная номенклатура активного сетевого, вычислительного и инфраструктурного оборудования, выбранного для реализации структурированной кабельной системы и локальной вычислительной сети этажа лаборатории. Все приведенные позиции активного телекоммуникационного оборудования и вычислительной техники соответствуют требованиям импортозамещения и присутствуют в Едином реестре российской радиоэлектронной продукции Минпромторга РФ.

\begin{small}
\begin{longtable}{|c|p{4.6cm}|p{3.1cm}|c|c|p{3.4cm}|}
\caption{Ведомость покупных изделий} \\
\hline
\textbf{№} & \textbf{Наименование и техническая} \newline \textbf{характеристика} & \textbf{Производитель /} \newline \textbf{Бренд} & \textbf{Ед.} \newline \textbf{изм.} & \textbf{Кол-} \newline \textbf{во} & \textbf{Примечание} \newline \textbf{(Назначение)} \\
\hline
\endfirsthead
\multicolumn{6}{c}%
{\tablename\ \thetable\ -- \textit{Продолжение на следующей странице}} \\
\hline
\textbf{№} & \textbf{Наименование и техническая} \newline \textbf{характеристика} & \textbf{Производитель /} \newline \textbf{Бренд} & \textbf{Ед.} \newline \textbf{изм.} & \textbf{Кол-} \newline \textbf{во} & \textbf{Примечание} \newline \textbf{(Назначение)} \\
\hline
\endhead
\hline \multicolumn{6}{r}{\textit{Продолжение на следующей странице}} \\
\endfoot
\hline
\endlastfoot

\multicolumn{6}{|l|}{\textbf{1 Активное сетевое оборудование}} \\
\hline
1.1 & Сервисный маршрутизатор ESR-1200 (1U, 10G SFP+) & Eltex & шт. & 1 & Серверная (пом. №\,19), шлюз сети \\
\hline
1.2 & Коммутатор ядра MES5324 (24x10G SFP+, 2x100G QSFP28) & Eltex & шт. & 1 & Серверная (пом. №\,19), уровень ядра \\
\hline
1.3 & Коммутатор доступа AQ-N3000-48P4Y2Q (48x1G PoE++) & Aquarius & шт. & 1 & Кроссовая (пом. №\,16), питание моноблоков \\
\hline
1.4 & Коммутатор доступа MES2448P (48x1G PoE+, 4x10G SFP+) & Eltex & шт. & 1 & Кроссовая (пом. №\,16), питание камер и СКУД \\
\hline
1.5 & Коммутатор доступа MES2448B (48x1G, 4x10G SFP+) & Eltex & шт. & 2 & Кроссовая (пом. №\,16), стенды лабораторий \\
\hline
\multicolumn{6}{|l|}{\textbf{2 Вычислительное оборудование и рабочие станции}} \\
\hline
2.1 & Медиасервер стоечный С2242И (2U) & Гравитон & шт. & 1 & Серверная, обработка видеотрафика \\
\hline
2.2 & Сервер автоматизированной сборки CI/CD С2122АА (2U) & Гравитон & шт. & 1 & Серверная, компиляция и прошивка \\
\hline
2.3 & Сетевое хранилище данных (NAS) серии QSRV (2U) & Qtech & шт. & 1 & Серверная, кэширование и архивы \\
\hline
2.4 & Моноблок оператора М75И (потребление до 120 Вт) & Гравитон & шт. & 6 & Пом. №\,20, рабочие места операторов \\
\hline
\multicolumn{6}{|l|}{\textbf{3 Системы безопасности и контроля (СКУД и СВН)}} \\
\hline
3.1 & Узконаправленная IP-камера QVC-MiR501Z1 (PoE+, 8 Вт) & Qtech & шт. & 6 & Контроль испытательных стендов \\
\hline
3.2 & Обзорная IP-камера QVC-MiR803M (PoE+, 7 Вт) & Qtech & шт. & 2 & Общий контроль помещений лабораторий \\
\hline
3.3 & Сетевой контроллер СКУД E4 PRO (PoE+, до 4,5 Вт) & Sigur & шт. & 2 & Двери помещений №\,4 и №\,11 \\
\hline
\multicolumn{6}{|l|}{\textbf{4 Системы бесперебойного электропитания (ИБП)}} \\
\hline
4.1 & ИБП стоечный РСК-Эксперт-II-6000С-М (6000 ВА / 6000 Вт) & Сайбер Электро & шт. & 1 & Серверная (пом. №\,19) \\
\hline
4.2 & ИБП стоечный РСК-Эксперт-II-6000 (6000 ВА) & Сайбер Электро & шт. & 1 & Кроссовая (пом. №\,16) \\
\hline
4.3 & Напольный батарейный кабинет БШ-20 (укомплектован АКБ 55 Ач) & Сайбер Электро & шт. & 1 & Для ИБП серверной комнаты \\
\hline
4.4 & Дополнительный стоечный батарейный модуль (3U) & Сайбер Электро & шт. & 1 & Для ИБП кроссовой комнаты \\
\hline
\multicolumn{6}{|l|}{\textbf{5 Инфраструктура и климат-контроль}} \\
\hline
5.1 & Шкаф телекоммуникационный напольный C3.RK426010.2P2P (42U) & C3 Solutions & шт. & 1 & Серверная комната \\
\hline
5.2 & Шкаф телекоммуникационный напольный C3.RI3303 (33U) & C3 Solutions & шт. & 1 & Кроссовая комната \\
\hline
5.3 & Прецизионный внутристоечный кондиционер C3.AIR.BIS007 (7,2 кВт) & C3 Solutions & шт. & 2 & Охлаждение шкафа серверной (схема N+1) \\
\hline
\end{longtable}
\end{small}

\textit{(Примечание: Раздел пассивного оборудования --- медный и оптический кабель, патч-панели, кабельные органайзеры, лотки и телекоммуникационные розетки --- будет добавлен в данную ведомость на заключительном этапе проектирования после проведения точного расчета контрольных длин кабельных трасс и выбора кабеленесущих систем).}

\end{appendices}

\end{document}
